\documentclass[10pt,a4paper]{amsart}
\usepackage[margin=19mm]{geometry}
\usepackage{amsmath,amssymb,mathtools}
\usepackage[scr=rsfs,cal=euler]{mathalfa}
\usepackage{xfrac}
\usepackage[unicode,hidelinks,bookmarks=false]{hyperref}
\newcommand{\ie}{i.e.\ }
\newcommand{\cf}{cf.\ }
\newcommand{\resp}{respectively\ }
\newcommand{\Subsection}{\subsection}
\providecommand{\nobreakdash}{\nobreak\hbox{-}}
\setlength{\emergencystretch}{2em}
\multlinegap0pt
\newtheorem{cor}{Corollary}
\newtheorem{lem}{Lemma}
\newtheorem{thm}{Theorem}
\title{On the Diophantine problem related to power circuits}
\author{Alexander Rybalov}
\date{Source: arXiv:2601.00835v4 (CC BY 4.0)}
\begin{document}
\maketitle

\noindent\emph{Source and licence.} On the Diophantine problem related to power circuits. Version arXiv:2601.00835v4 (CC BY 4.0), distributed by arXiv under the Creative Commons Attribution 4.0 International licence, http://creativecommons.org/licenses/by/4.0/, as recorded in the arXiv licence metadata for this exact version and shown on the abstract page https://arxiv.org/abs/2601.00835v4. Full licence notice: \texttt{licenses/arxiv-2601.00835-CC-BY-4.0.txt}.\par\smallskip

\noindent\textit{Editorial scope.} Counted unit: the article's thm main (source0.tex lines 269--271). The article's complete original proof of it is delivered with it (source0.tex lines 272--292, one \texttt{proof} environment of 21 lines). No mathematics is written, repaired or re-derived: the statement and the proof are the article's own lines, in the article's own notation, and no retained line is substituted or re-flowed.  Retained uncounted as the source-local context the proof needs: lines 137--139; lines 192--194; lines 257--259.  Nothing was omitted from any retained span, so the package carries no omission record.  No retained line contains a citation, so the wrapper carries no bibliography and no bibliography entry.  No retained line contains a figure environment, an image-inclusion command or drawing code, so no image is delivered and the package declares no asset dependency.  Editorial formatting only, in addition: an AMS \texttt{amsart} preamble that restates the article's own declarations and macros used by the retained lines, namely the environments \texttt{cor}, \texttt{lem}, \texttt{thm} on separate counters, together with the standard packages those lines need.  The retained environments print in this document as Corollary 1, Lemma 1, Theorem 1 in order of appearance, whereas the article prints the same items with the numbering of its own full document.  The complete native source of the article as distributed in the arXiv e-print for this exact version is bundled unchanged as \texttt{sources/192033/source0.tex}.
\begin{cor} \label{dpk_undec}
For every natural number $k$ the problem $\mathcal{DP}(\mathbb{N}_{>k})$ is undecidable.
\end{cor}

\begin{lem} \label{str_def}
The strict order relation $x < y$ is Diophantine definable in $\widetilde{N}$.
\end{lem}

\begin{lem} \label{mult_def}
The multiplication operation $mul(x,y) = xy$ for $x,y>1$ is Diophantine definable in $\widetilde{N}$.
\end{lem}

\begin{thm} \label{main}
The Diophantine problem over $\widetilde{N} = \langle \mathbb{N}_{>0}; +, x \cdot 2^y, \leq, 1 \rangle$ is undecidable.
\end{thm}

\begin{proof}
We reduce $\mathcal{DP}(\mathbb{N}_{>1})$ to the Diophantine problem over $\widetilde{N}$ in the following way.
Without loss of generality any given nontrivial system $S$ of Diophantine equations over $\mathbb{N}_{>1}$ 
consists of equations of the form $P(x_1, \ldots, x_n) = Q(x_1, \ldots, x_n)$,
where $P$ and $Q$ are non-zero polynomials with positive integer coefficients.
Such system $S$ can be transformed to an equivalent system in the {\it Skolem form}, consisting of equations of
the following types:
\begin{enumerate}
\item $x_i = x_j x_k$,
\item $x_i = x_j + x_k$,
\item $x_i = x_j + 1$,
\item $x_i = x_j$.
\end{enumerate}
By Lemma \ref{mult_def}, we can replace every equation of type (1) by an equivalent system of equations over $\widetilde{N}$.
Also for every variable $x$, which is included in equations of types (2), (3) or (4), but not included in any equation of type (1),
we add the Diophantine condition $x>1$ using Lemma \ref{str_def}. 
By construction, the obtained system of equations over $\widetilde{N}$ is equivalent to the original system $S$ over $\mathbb{N}_{>1}$.

By Corollary \ref{dpk_undec}, the problem $\mathcal{DP}(\mathbb{N}_{>1})$ is undecidable;
therefore, $\mathcal{DP}(\widetilde{N})$ is also undecidable.
\end{proof}

\end{document}
